{\Large\bfseries Mockup\par}
\bigskip

\fullname

30 minutes, one handwritten or 4x6 index card (or A6) cheat sheet allowed, no printed material.

\bigskip
\begin{center}\rule{0.4\linewidth}{0.4pt}\end{center}
\bigskip

\begin{question}
In the statements below, $N$ is the size of the input, complexities are worst-case, and all algorithms solve the same problem and run on the same computer.

Exactly \textbf{2} of the following statements are true. Select them.
\begin{itemize}[leftmargin=2em,itemsep=1pt,topsep=2pt]
\item[\cb{1}] an $O(1)$ algorithm can take longer than an $O(N)$ algorithm on an input of size $N = 10$
\item[\cb{0}] a constant-time algorithm does the same number of operations for every $N$
\item[\cb{0}] an $O(\log N)$ algorithm is faster than an $O(N)$ algorithm for every $N$
\item[\cb{0}] knowing only that one algorithm is $O(N)$ and another is $O(N^2)$, I can compute the $N$ above which the first one is faster
\item[\cb{1}] for $N$ large enough, an $O(N \log N)$ algorithm is faster than an $O(N^2)$ algorithm
\end{itemize}
\end{question}

\begin{question}
\lstinline|list| is an empty list, which can be either a \lstinline|java.util.ArrayList| or a \lstinline|java.util.LinkedList|, and \lstinline|N| is a positive integer.

\begin{lstlisting}
for (int i = 0; i < N; i++)
    list.add(i);
while (!list.isEmpty())
    System.out.println(list.remove(0));
\end{lstlisting}

Which implementation gives the lowest worst-case number of operations for this code (select one)? Then write that complexity as a function of $N$.

\begin{itemize}[leftmargin=2em,itemsep=1pt,topsep=2pt]
\item[\cb{0}] \lstinline|ArrayList|
\item[\cb{1}] \lstinline|LinkedList|
\item[\cb{0}] both have the same complexity
\end{itemize}

Complexity: $O($\underline{\hspace{1in}}$)$ \ifanswers ($N$)\fi
\end{question}

\begin{question}
In a singly linked list (nodes with fields \lstinline|data| and \lstinline|next|), \lstinline|before| refers to a node that is not the last one.

The four lines below remove the node after \lstinline|before| from the list, detach it by setting its \lstinline|next| to \lstinline|null|, and return its data, but they are in the wrong order.
Write the numbers 1 to 4 in the boxes to give the order in which they must be executed.

\medskip
\begin{tabular}{@{}l@{\quad}l@{}}
\ansbox{4} & \lstinline|return gone.data;| \\[3pt]
\ansbox{3} & \lstinline|gone.next = null;| \\[3pt]
\ansbox{1} & \lstinline|Node gone = before.next;| \\[3pt]
\ansbox{2} & \lstinline|before.next = gone.next;| \\
\end{tabular}
\end{question}

\begin{question}
Consider the usual implementation of a singly-linked list \lstinline|Node|.

\begin{lstlisting}
class Node {
    int data;
    Node next;
    Node(int data) { this.data = data; }
}
\end{lstlisting}

What does the following code fragment print?

Hint: draw the nodes and the references on scrap paper, and update the drawing after each line.

\begin{lstlisting}
Node head = new Node(7);
head.next = new Node(8);
head.next.next = new Node(9);
Node p = head;
Node q = head.next.next;
q.next = p;
head = p.next;
p.next = null;
q.data = 6;
for (Node cur = head; cur != null; cur = cur.next)
    System.out.print(cur.data + " ");
\end{lstlisting}

Answer: \underline{\hspace{2in}} \ifanswers (8 6 7)\fi
\end{question}

\begin{question}
The following method of a doubly linked list removes the node \lstinline|n| from the list.
\lstinline|head| and \lstinline|tail| refer to the first and the last node (\lstinline|null| if the list is empty), and \lstinline|size| is the number of elements.
Complete the missing line, so that the method works correctly.

\begin{lstlisting}
private double unlink(Node n) {
    if (n.prev == null) head = n.next;
    else n.prev.next = n.next;
    if (n.next == null) ___
    else n.next.prev = n.prev;
    size--;
    return n.data;
}
\end{lstlisting}

Answer:
\begin{itemize}[leftmargin=2em,itemsep=1pt,topsep=2pt]
\item[\cb{0}] \lstinline|tail = n;|
\item[\cb{0}] \lstinline|tail = n.next;|
\item[\cb{1}] \lstinline|tail = n.prev;|
\item[\cb{0}] \lstinline|tail.prev = n.prev;|
\item[\cb{0}] \lstinline|n.prev = tail;|
\end{itemize}
\end{question}

\begin{question}
The following method of an array-based list of \lstinline|int|s should insert \lstinline|e| at position \lstinline|pos|, but one line contains a bug.
\lstinline|data| is the array holding the elements and \lstinline|size| is the number of elements in the list.
Which change fixes the method (select one)?

Hint: you can simulate by hand what the method does on a very small list to pin down which operation is buggy.

\begin{lstlisting}[numbers=left,numberstyle=\tiny,xleftmargin=2em]
public void add(int pos, int e) {
    if (pos < 0 || pos > size) throw new IndexOutOfBoundsException();
    if (size == data.length) grow();
    for (int i = pos; i < size; i++)
        data[i + 1] = data[i];
    data[pos] = e;
    size++;
}
\end{lstlisting}

\begin{itemize}[leftmargin=2em,itemsep=1pt,topsep=2pt]
\item[\cb{0}] line 2 becomes \lstinline!if (pos < 0 || pos >= size) throw new IndexOutOfBoundsException();!
\item[\cb{0}] line 3 becomes \lstinline|if (size == data.length - 1) grow();|
\item[\cb{1}] line 4 becomes \lstinline|for (int i = size - 1; i >= pos; i--)|
\item[\cb{0}] line 5 becomes \lstinline|data[i] = data[i + 1];|
\item[\cb{0}] line 6 becomes \lstinline|data[pos + 1] = e;|
\end{itemize}
\end{question}

\begin{question}
\lstinline|list| is a \lstinline|java.util.ArrayList| or a \lstinline|java.util.LinkedList| containing $N$ elements.
Costs are worst-case (amortized for adding at the end of an \lstinline|ArrayList|).

Exactly \textbf{2} of the following operations are $O(1)$ on an \lstinline|ArrayList| but $O(N)$ on a \lstinline|LinkedList|. Select them.
\begin{itemize}[leftmargin=2em,itemsep=1pt,topsep=2pt]
\item[\cb{1}] read the element at a given position
\item[\cb{0}] add an element at the beginning
\item[\cb{0}] add an element at the end
\item[\cb{1}] replace the element at a given position
\item[\cb{0}] check if the list contains a given element
\end{itemize}
\end{question}

\begin{question}
\lstinline|list| contains $N$ elements, and the following method is called on it.

\begin{lstlisting}
int process(List<Integer> list) {
    int negatives = 0;
    for (int x : list)
        if (x < 0) negatives++;
    for (int i = 0; i < 10; i++)
        list.add(0, i);
    return negatives;
}
\end{lstlisting}

What is the worst-case complexity of \lstinline|process|, as a function of $N$, when \lstinline|list| is a \lstinline|java.util.ArrayList|? And when it is a \lstinline|java.util.LinkedList|?

\lstinline|ArrayList|: \underline{\hspace{1in}} \ifanswers ($O(N)$: the loop is $O(N)$ and the 10 calls to \lstinline|add(0, i)| cost $O(N)$ each)\fi \qquad
\lstinline|LinkedList|: \underline{\hspace{1in}} \ifanswers ($O(N)$: the loop is $O(N)$ and the 10 calls to \lstinline|add(0, i)| cost $O(1)$ each)\fi
\end{question}

\begin{question}
The following inner class of a doubly linked list of \lstinline|double|s (fields \lstinline|head|, \lstinline|tail|, \lstinline|size|) is an iterator over its elements, from last to first.
Write the last three lines of \lstinline|next()|, so that the iterator returns every element of the list exactly once, from last to first.

\begin{lstlisting}
private class DescendingItr implements Iterator<Double> {
    private Node current = tail;    // node holding the next element to return

    public boolean hasNext() { return current != null; }

    public Double next() {
        if (!hasNext()) throw new NoSuchElementException();
        // write the three lines in the space below this block
        // write the three lines in the space below this block
        // write the three lines in the space below this block
    }
}
\end{lstlisting}

\ifanswers
\begin{lstlisting}
        double e = current.data;
        current = current.prev;
        return e;
\end{lstlisting}
\else
\vspace{1.2in}
\fi
\end{question}

\begin{question}
(2 points) The list \lstinline|[1, 2, 3]| is represented as a singly linked list of \lstinline|Integer|s \textbf{with a sentinel} node (a first node that holds no element).
Draw a picture of its representation in memory (including the \lstinline|LinkedList| object, all its \lstinline|Node|s, and the data contained by the nodes), with arrows representing references, including \lstinline|null| references.

Hint: remember that the \lstinline|Integer| class is not a primitive type; it is an object that stores a single \lstinline|int| value.

\ifanswers
\begin{verbatim}
LinkedList: head -> S, tail -> N3, size = 3

S:  data null,          next -> N1   (sentinel)
N1: data -> Integer(1), next -> N2
N2: data -> Integer(2), next -> N3
N3: data -> Integer(3), next null

3 distinct Integer objects
\end{verbatim}
\else
\vspace{3in}
\fi
\end{question}
